\subsubsection*{2.6.}
设 $\widetilde D^*$ 是 $D^*$ 的一个万有覆盖，$k(\bar\eta)$ 是 $k(\eta)$ 的相应代数闭包，并设 $\overline S$ 是 $S^h$ 在 $k(\bar\eta)$ 中的正规化。将 $\overline S$ 写成仿射 $S$-概形 $S_i$ 的逆极限。对每个 $i$，将 $D$ 换为较小圆盘 $D_i$ 后，有自然的 $S_{\mathrm{cl}}$-空间态射
\[
a_i:\widetilde D_i\longrightarrow(S_i)_{\mathrm{cl}}.
\]
设 $j_i:U_i\hookrightarrow S_i$ 是 $S_i$ 中 $\overline S$ 的闭点的像 $0$ 的开补集，并令 $i:0\hookrightarrow S_i$。

设 $F$ 是 $X$ 上的层，并令 $F_{\mathrm{cl}}=\varepsilon^*F$。仍用 $j_i$ 表示 $X\times_SU_i$ 到 $X\times_SS_i$ 的包含，用 $i$ 表示 $X_0$ 到 $X\times_SS_i$ 的包含，用 $F$ 表示 $F$ 在 $X\times_SS_i$ 上的逆像。则
\[
\Psi(F)=
\left(i^*F,\ \varinjlim_i i^*j_{i*}j_i^*F,\
\varinjlim_i\varphi_i\right).
\]

函子 $\varepsilon^*$ 与逆像交换。因此态射
\[
\varepsilon^*j_{i*}\longrightarrow(j_{i,\mathrm{cl}})_*\varepsilon^*
\]
给出
\[
c_i:
(i^*j_{i*}j_i^*F)_{\mathrm{cl}}
\longrightarrow i^*j_{i*}j_i^*(F_{\mathrm{cl}}).
\]

设 $\overline j$（分别为 $\overline j_i$）表示 $\widetilde D^*$（分别为 $\widetilde D_i^*$）到 $\widetilde D$（分别为 $\widetilde D_i$）的包含，也表示所诱导的包含
\[
X\times_S\widetilde D^*\hookrightarrow X\times_S\widetilde D,
\qquad
X\times_S\widetilde D_i^*\hookrightarrow X\times_S\widetilde D_i.
\]
伴随映射 $\mathrm{Id}\longrightarrow a_{i*}a_i^*$ 定义
\[
i^*j_{i*}j_i^*F_{\mathrm{cl}}
\longrightarrow
i^*\overline j_{i*}\overline j_i^*\overline F_{\mathrm{cl}}
=i^*\overline j_*\overline j^*\overline F_{\mathrm{cl}}.
\]
取极限得到
\[
c_\eta:\Psi_{\bar\eta}(F)_{\mathrm{cl}}
\longrightarrow\bigl(\Psi_{\mathrm{cl}}(F_{\mathrm{cl}})\bigr)_{\bar\eta}.
\EditorialFootnote{这一未编号公式的源是固定万有覆盖所对应的几何泛分量；目标也应取同一分量。原文目标未标出该分量。此处补出分量记号，与 1.2.4(c) 的函子描述及后面的 (2.6.2)、2.7 相容。}
\]
此态射定义\underline{比较态射}
\begin{equation}\tag{2.6.1}
c:\varepsilon^*\Psi(F)
\longrightarrow\Psi_{\mathrm{cl}}(\varepsilon^*F).
\end{equation}
同样定义
\begin{equation}\tag{2.6.2}
c_\eta:\varepsilon^*\Psi_\eta(F)
\longrightarrow\Psi_\eta(\varepsilon^*F).
\end{equation}

